% !TEX root = ../raiz.tex

\chapter{Revisão teórica}

\section{Princípios físicos da TIE}
\label{subsec:principios_fisicos_tie}

O princípio físico da TIE está relacionado ao fato de que tecidos biológicos distintos apresentam propriedades elétricas diferentes. Regiões com maior conteúdo de água e eletrólitos tendem a apresentar menor impedância, enquanto estruturas como gordura, osso e ar apresentam comportamento elétrico distinto, o que permite inferir alterações internas a partir das medidas realizadas externamente \cite{Lian2025}. No tórax, essa característica é particularmente relevante porque a ventilação pulmonar modifica a impedância regional ao longo do ciclo respiratório, tornando a TIE adequada para o acompanhamento dinâmico da distribuição de ar nos pulmões. Em termos práticos, o aumento da fração de ar nos pulmões durante a inspiração está associado ao aumento da impedância elétrica regional.

Do ponto de vista matemático, a técnica envolve um problema inverso não linear e mal posto, pois a distribuição interna de condutividade precisa ser estimada a partir de tensões medidas apenas na superfície. Essa característica limita a resolução espacial das imagens reconstruídas quando comparadas à tomografia computadorizada e à ressonância magnética, embora a TIE apresente vantagens importantes em resolução temporal, custo e possibilidade de monitoramento contínuo \cite{KhanLing2019}.


\section{Aplicações clínicas e potencial como técnica complementar}
\label{subsec:aplicacoes_tie}

Na prática clínica, a aplicação mais consolidada da TIE está no monitoramento pulmonar à beira do leito, especialmente em pacientes críticos submetidos à ventilação mecânica. Nesse contexto, a técnica pode ser utilizada para acompanhar a distribuição regional da ventilação e da perfusão, auxiliando na avaliação de colapso alveolar, hiperdistensão, recrutamento pulmonar e resposta a ajustes ventilatórios \cite{Scaramuzzo2024}. Esse perfil torna a TIE especialmente útil em situações nas quais o acompanhamento contínuo da função pulmonar é mais importante do que a obtenção de uma imagem anatômica de alta resolução.

No contexto brasileiro, há relatos de uso clínico da TIE para avaliação da ventilação pulmonar regional em ambiente hospitalar. Um exemplo é o estudo publicado no \textit{Jornal Brasileiro de Pneumologia}, no qual a técnica foi utilizada para avaliar uma paciente com estenose brônquica unilateral pós-tuberculose, produzindo resultados semelhantes aos da cintilografia de ventilação/perfusão, com a vantagem de fornecer avaliação dinâmica sem exposição à radiação \cite{Marinho2013JBP}. Embora não substitua métodos anatômicos de alta resolução, a TIE apresenta utilidade particular no monitoramento funcional contínuo, sobretudo em aplicações torácicas.

\section{Formulação matemática}
\label{sec:fundamentos_tie}

A Tomografia por Impedância Elétrica (TIE) baseia-se na aplicação de correntes elétricas alternadas de pequena amplitude através de eletrodos posicionados na fronteira de um domínio (ou região de interesse) condutivo e na medição das tensões elétricas resultantes nos demais eletrodos. No contexto adotado neste trabalho, considera-se uma seção transversal do domínio, com os eletrodos distribuídos em anel ao longo de sua fronteira. A cada etapa da aquisição, um par de eletrodos é utilizado para a injeção de corrente, enquanto as tensões são medidas nos pares restantes segundo um protocolo previamente definido. Em seguida, o par de injeção é alterado, e o processo é repetido até que todos os padrões de excitação previstos sejam completados. Um ciclo completo dessas excitações corresponde a um \textit{frame} de aquisição.

Embora o nome da técnica faça referência à impedância, os métodos utilizados neste trabalho reconstruem a distribuição de condutividade elétrica, denotada por \(\sigma\). A terminologia ``impedância'' decorre do uso de corrente alternada no processo de medição, mas a formulação dos algoritmos de reconstrução empregados é expressa em termos de condutividade. Essa distinção é importante para manter consistência com a representação adotada na aplicação, na qual as imagens reconstruídas são interpretadas como mapas de variação de condutividade, \(\Delta \sigma\).

\subsection{Problema direto}
\label{subsec:problema_direto}

O problema direto da TIE consiste em determinar quais tensões seriam medidas nos eletrodos quando a distribuição de condutividade da região de interesse é conhecida. Em regime quase-estático e sob hipóteses usuais de condução elétrica em meios biológicos, o potencial elétrico \(u\) no interior do domínio \(\Omega\) é descrito pela equação

\begin{equation}
\nabla \cdot (\sigma \nabla u) = 0 \quad \text{em } \Omega ,
\end{equation}

em que \(\sigma\) representa a distribuição espacial de condutividade e \(u\) o potencial elétrico. As condições de contorno associadas a essa equação modelam a injeção de corrente e a medição de tensões nos eletrodos. Em formulações mais realistas, essa interação é frequentemente representada pelo modelo completo de eletrodos (\textit{Complete Electrode Model} -- CEM), que incorpora explicitamente a interface entre os eletrodos e o meio condutivo \cite{Somersalo1992CEM}.

No presente trabalho, o problema direto foi resolvido numericamente por meio do Método dos Elementos Finitos (FEM), utilizando a implementação disponível no framework pyEIT \cite{pyEITReadme}. Para isso, o domínio foi representado por uma malha triangular, que constitui a discretização numérica da região de interesse. Essa malha é a mesma estrutura posteriormente utilizada no problema inverso, o que faz com que sua configuração influencie diretamente a reconstrução final. Em termos operacionais, os parâmetros de malha ajustáveis na interface desenvolvida, como geometria e densidade, afetam não apenas a visualização, mas também a forma como o problema direto é resolvido internamente.

\subsection{Problema inverso e caráter mal posto}
\label{subsec:problema_inverso}

O problema inverso da TIE consiste em estimar a distribuição de condutividade no interior do domínio a partir das tensões medidas nos eletrodos. Diferentemente do problema direto, essa etapa não é, em geral, bem posta no sentido de Hadamard. Em particular, a solução pode não ser única, pode não existir para dados ruidosos e, sobretudo, pode não depender de forma estável das medições. Em consequência, pequenas perturbações nos dados de entrada podem produzir grandes variações na imagem reconstruída \cite{Lionheart2003EITRA}.

Esse comportamento está diretamente relacionado à natureza limitada das medições de fronteira. O número de medições independentes obtidas em cada \textit{frame} é muito menor do que o número de incógnitas associado à discretização espacial do domínio, o que faz com que a reconstrução dependa de restrições adicionais, aproximações ou estratégias de regularização para produzir soluções estáveis e interpretáveis. Esse é um dos principais motivos pelos quais diferentes algoritmos de reconstrução podem gerar imagens com características bastante distintas a partir do mesmo conjunto de dados experimentais.

\subsection{Reconstrução por diferença temporal}
\label{subsec:diferenca_temporal}

A interface desenvolvida utilizou reconstrução por diferença temporal (\textit{time-difference EIT}), em vez de reconstrução absoluta. Nessa abordagem, não se busca estimar diretamente a distribuição absoluta de condutividade do domínio, mas sim a variação de condutividade em relação a um estado de referência. Em outras palavras, cada imagem reconstruída representa \(\Delta \sigma\), isto é, a mudança de condutividade entre um \textit{frame} de medição e um \textit{frame} de referência previamente definido \cite{Putensen2019}.

No sistema implementado, o primeiro \textit{frame} do conjunto de dados foi utilizado como referência, sendo armazenado internamente antes do início da reconstrução dos demais quadros. No código desenvolvido, essa etapa é realizada pelo método \texttt{setVref} da classe \texttt{EITsolver}, que armazena as medições de referência antes do início da reconstrução sequencial dos demais \textit{frames}. Essa abordagem é compatível com aplicações práticas da TIE, nas quais a reconstrução por diferença temporal tende a ser mais robusta a imperfeições do modelo, incertezas geométricas e variações no contato eletrodo-superfície \cite{Putensen2019}. Consequentemente, as imagens exibidas pela interface representam mudanças relativas de condutividade ao longo do tempo, e não a condutividade absoluta do meio.

\subsection{Métodos de reconstrução de imagens}
\label{subsec:metodos_reconstrucao_imagens}

Foram incorporados três métodos de reconstrução presentes no ecossistema do pyEIT: Back-Projection (BP), Jacobiano (JAC) e GREIT. A documentação do framework descreve esses algoritmos como parte do conjunto principal de métodos inversos disponibilizados pela biblioteca, juntamente com o solver direto baseado em elementos finitos \cite{pyEITReadme}. Embora todos operem sobre o mesmo conjunto de medições elétricas de fronteira, cada um aborda o problema inverso de maneira distinta, o que resulta em diferenças de contraste, definição espacial, estabilidade temporal e custo computacional.

\subsubsection{Back-Projection (BP)}
\label{subsubsec:bp}

O método Back-Projection constitui a abordagem conceitualmente mais simples entre as empregadas neste trabalho. Em termos gerais, cada diferença de tensão medida é retroprojetada sobre o domínio segundo a região de sensibilidade associada ao respectivo par de medição, e a imagem final é obtida pela superposição dessas contribuições. Essas regiões de sensibilidade são definidas pela matriz Jacobiana, que relaciona cada medição de tensão às variações de condutividade em cada elemento da malha, ainda que o BP não realize a inversão explícita dessa matriz. No ecossistema do pyEIT, o solver correspondente é descrito como um método de \textit{back projection} e \textit{filtered back projection} \cite{pyEITReadme}.

Do ponto de vista computacional, trata-se de um método leve, com baixa complexidade e sem a necessidade de resolver iterativamente um sistema inverso regularizado. Essa simplicidade torna o BP atrativo quando se deseja reconstrução rápida ou atualização imediata das imagens. Por outro lado, trabalhos dedicados a esse tipo de abordagem destacam que o algoritmo clássico de retroprojeção em TIE apresenta baixa resolução espacial e sensibilidade não uniforme ao longo do domínio, resultando em imagens com maior borramento e presença de artefatos \cite{BarberBrown1984}.

Nos resultados obtidos neste trabalho, o BP permitiu localizar qualitativamente a posição da haste no interior do meio, mas com baixa definição de contorno e presença de regiões difusas que extrapolavam a área fisicamente perturbada. Também foram observados artefatos de sinal oposto em regiões sem correspondência física aparente, o que é coerente com a maior sensibilidade do método ao ruído e à ausência de mecanismos explícitos de regularização.

Para atenuar oscilações entre \textit{frames} consecutivos, foi implementado especificamente para esse método um parâmetro de suavização temporal \(\alpha\), aplicado segundo a expressão

\begin{equation}
\label{eq:suavizacao_temporal}
\hat{\sigma}_k = \alpha \cdot \hat{\sigma}_{k-1} + (1 - \alpha)\cdot \sigma_k ,
\end{equation}

em que \(\hat{\sigma}_k\) representa a imagem suavizada no quadro \(k\), \(\hat{\sigma}_{k-1}\) a imagem suavizada do quadro anterior e \(\sigma_k\) a reconstrução bruta obtida no quadro corrente. Quando \(\alpha = 0\), não há suavização temporal e a imagem exibida corresponde diretamente à reconstrução bruta; à medida que \(\alpha\) se aproxima de 1, a imagem passa a variar mais lentamente no tempo, reduzindo a cintilação visual, mas também aumentando a inércia da resposta.

\subsubsection{Jacobiano (JAC)}
\label{subsubsec:jac}

O método Jacobiano aborda o problema inverso a partir da matriz de sensibilidade \(J\), cujos elementos relacionam variações locais de condutividade às variações observadas nas tensões medidas. Nessa formulação, o vetor de diferenças de tensão \(\Delta v\) é aproximado por uma relação linear com a variação de condutividade \(\Delta \sigma\), permitindo escrever a reconstrução como um problema inverso regularizado. No ecossistema do pyEIT, esse solver é descrito como parte da família tradicional de métodos de Gauss--Newton, com suporte a diferentes estratégias internas de regularização, incluindo \texttt{lm} e \texttt{kotre} \cite{pyEITReadme}.

Em sua forma linearizada, a solução pode ser expressa como

\begin{equation}
\Delta \sigma = (J^T J + \lambda R)^{-1} J^T \Delta v ,
\end{equation}

em que \(\lambda\) representa o parâmetro de regularização e \(R\) a matriz regularizadora. O parâmetro \(\lambda\) controla o compromisso entre aderência aos dados medidos e estabilidade da solução: valores maiores tendem a produzir imagens mais suaves e menos sensíveis a ruído, enquanto valores menores favorecem maior contraste, mas também maior instabilidade. Na implementação utilizada neste trabalho, a opção \texttt{kotre} aplica uma regularização ponderada pela distribuição de sensibilidade da malha, enquanto a opção \texttt{lm} utiliza uma forma mais uniforme baseada nos termos diagonais de \(J^T J\) \cite{pyEITReadme}.

Em comparação com o BP, o método Jacobiano depende explicitamente de um modelo linearizado do problema inverso e da escolha de parâmetros de regularização. Esse tipo de formulação é característico dos métodos regulares de reconstrução em TIE, nos quais a estabilidade da solução depende diretamente das restrições impostas ao problema mal posto \cite{AdlerHolder2021}. Como consequência, pequenas mudanças nos parâmetros podem alterar de forma perceptível o aspecto final da imagem reconstruída.

Nos dados experimentais utilizados neste trabalho, o JAC produziu imagens com maior definição espacial do que o BP, com contraste mais localizado e regiões de variação de condutividade visualmente mais nítidas. Ainda assim, também apresentou sensibilidade perceptível à escolha da regularização, o que se refletiu em diferenças marcantes entre \textit{frames} subsequentes em algumas sequências analisadas. Em especial, os padrões bipolares mais intensos observados em determinados \textit{frames} e a maior instabilidade entre algumas reconstruções consecutivas são compatíveis com a própria natureza do método, e não com uma limitação da interface desenvolvida.

\subsubsection{GREIT}
\label{subsubsec:greit}

O método GREIT (\textit{Graz consensus Reconstruction algorithm for Electrical Impedance Tomography}) adota uma estratégia distinta das duas anteriores. Em vez de resolver diretamente o problema inverso a cada novo conjunto de medições, o método constrói previamente uma matriz de reconstrução linear durante a etapa de configuração. A partir dessa matriz, cada nova imagem é obtida por uma operação linear direta sobre o vetor de diferenças de tensão \(\Delta v\), o que torna a reconstrução computacionalmente eficiente após a fase inicial de preparação \cite{GREIT2009}.

A formulação original do GREIT foi proposta como uma abordagem unificada para reconstrução bidimensional em TIE pulmonar, estruturada a partir de modelos por elementos finitos, figuras de mérito padronizadas e um procedimento sistemático para obtenção de uma matriz de reconstrução otimizada. Entre essas figuras de mérito, destacam-se a resposta de amplitude uniforme, o erro de posição reduzido, a limitação de artefatos de \textit{ringing}, a resolução uniforme e a deformação de forma controlada \cite{GREIT2009}. No ecossistema do pyEIT, o GREIT é descrito como um algoritmo baseado em filtragem espacial estatística e em uma matriz de reconstrução de imagens, com saída em grade regular bidimensional \cite{pyEITReadme}.

Uma consequência prática dessa formulação é que, após a etapa de configuração, o processo de reconstrução se torna muito rápido, pois se reduz essencialmente a uma multiplicação matricial. Em contrapartida, o método passa a depender da estrutura previamente definida para essa matriz de reconstrução, tornando-se menos flexível do que abordagens baseadas diretamente na Jacobiana quando se deseja adaptar a formulação a diferentes hipóteses do problema. Além disso, como a saída do GREIT é produzida diretamente em uma grade regular, e não sobre a malha triangular do domínio, a visualização final pode apresentar limites retangulares característicos da grade utilizada \cite{pyEITReadme}.

Nos resultados obtidos na interface, o GREIT apresentou imagens visualmente mais suaves e estáveis entre \textit{frames} consecutivos, com localização consistente da haste em relação aos demais métodos. Embora o contraste visual tenha sido, em alguns casos, menor do que no JAC, a maior uniformidade espacial e temporal facilitou a leitura qualitativa do deslocamento do objeto. A presença visível de uma borda retangular na saída não foi interpretada como artefato físico do experimento, mas como consequência direta do formato de grade regular adotado internamente pelo método.

De modo geral, os três métodos implementados localizaram a haste na mesma região do domínio ao longo das reconstruções, o que sugere consistência qualitativa na informação espacial extraída das medições, mesmo diante do caráter mal posto do problema inverso. As diferenças entre eles manifestaram-se principalmente na definição dos contornos, no contraste, na estabilidade entre \textit{frames} subsequentes e no custo computacional associado a cada abordagem. Tornar essas diferenças visíveis e exploráveis foi precisamente um dos objetivos centrais da interface desenvolvida neste trabalho. 

\section{Padrão de Arquitetura MVC}
\label{sec:mvc}

O padrão de arquitetura MVC (\textit{Model--View--Controller}) foi proposto originalmente por Trygve Reenskaug em 1979, durante suas atividades de pesquisa no Xerox Palo Alto Research Center (PARC), com o objetivo de organizar aplicações interativas de forma mais estruturada e flexível \cite{Reenskaug1979}. A motivação central do padrão era separar explicitamente as responsabilidades relacionadas à manipulação de dados, à apresentação visual e à interação com o usuário, reduzindo o acoplamento entre essas partes e favorecendo a evolução incremental dos sistemas.

No MVC, o componente \textit{Model} é responsável por gerenciar os dados da aplicação e encapsular a lógica da aplicação associada, sem possuir qualquer conhecimento direto sobre a forma como essas informações são apresentadas ao usuário. A \textit{View} corresponde à camada de apresentação, encarregada de exibir visualmente os dados fornecidos pelo modelo, traduzindo estados internos em representações gráficas compreensíveis. O \textit{Controller}, por sua vez, atua como intermediário entre o usuário e o sistema, recebendo eventos de entrada, interpretando ações e coordenando as interações entre o modelo e a visualização.

A relação entre esses três componentes estabelece um fluxo bem definido de controle e dados. As ações do usuário são inicialmente capturadas pelo \textit{Controller}, que interpreta essas entradas e aciona as operações apropriadas no \textit{Model}. Após a atualização do estado interno, o modelo notifica a \textit{View}, que então se atualiza para refletir as novas informações, mantendo a consistência entre dados e apresentação \cite{GoF1994}. 
% FIGURA: diagrama MVC — a ser inserida posteriormente

A adoção do padrão MVC é particularmente comum em aplicações gráficas interativas devido à clara separação de responsabilidades que ele promove, facilitando a manutenção, a testabilidade e a extensão do software. Além disso, essa organização permite que diferentes camadas de visualização sejam substituídas ou modificadas sem impacto direto sobre a lógica central da aplicação. No presente trabalho, a aplicação EITduino adotou uma arquitetura inspirada nesse padrão, conforme descrito na Seção~\ref{sec:arquitetura_aplicacao}.